\hypertarget{native-291-remarks}{}\pdfbookmark[1]{Bemerkungen zu einzelnen Stellen}{bemerkungen-remarks}
Zu Seite 45, Zeile 1 findet sich eine Notiz von \textsc{Gauss} in seinem Handexemplar,
wonach dort $\omega'$ statt $\omega$ stehen soll. Allerdings ist $\omega$ richtig, falls in
Gleichung XI., Art. 22, nur das zweite Glied mit Hilfe hyperbolischer Logarithmen berechnet
wird, das erste ($e\operatorname{tang} F$) aber mit Briggs’schen Logarithmen. $\omega'$ muss stehen, wenn \textls[150]{beide}
Glieder mit hyperbolischen Logarithmen berechnet werden. Dann muss aber $\lambda$ im
Nenner fortfallen.

Zu den numerischen Beispielen ist im Einzelnen Folgendes zu bemerken. Auf den Planeten
\textls[150]{Juno} beziehen sich außer Art. 150---155 auch die Zahlen in den Art. 10, 13, 14, 51,
63---65, 69, 77, 87\,I, 97\,I und die Berechnung des Sonnenorts Art. 73. Auf den Planeten
\textls[150]{Pallas} außer Art. 156---157 auch Art. 56 und 58 und die Berechnung des Sonnenorts
Art. 70. Auf den Planeten \textls[150]{Ceres} außer Art. 159 auch Art. 87\,II, 97\,II. Außerdem stehen
im Zusammenhang die Beispiele in den Art. 38, 43, 97\,III für eine Ellipse mit großer
Exzentrizität, und Art. 23, 24, 26, 46, 105 für eine Hyperbel.

Zu Art. 150---155. Die Beobachtungen der Juno teilt \textsc{Maskelyne} in einem Brief an
\textsc{Gauss} vom 8. Mai 1805 mit. Die Sonnenörter, sowie die Nutation und die Schiefe der
Ekliptik hat \textsc{Gauss} aus den \textsc{Zach}schen Sonnentafeln entnommen. Auf eine
genaue Kontrolle derselben (Seite 195, Zeile 8---10) wurde verzichtet, da ihre strenge
Richtigkeit für die Durchführung der Beispiele unwesentlich ist und sie mit Rücksicht auf
die weitere Rechnung auch nicht hätten geändert werden können. Allerdings ist Seite 195, Zeile 8
im ersten Sonnenort $53''{,}72$ in $53''{,}71$ verwandelt worden, um das Folgende hiermit zur
Übereinstimmung zu bringen. Die Werte der Konstanten, welche \textsc{Gauss} zur Reduktion
der Beobachtungen benutzt hat, sind
\[
\begin{aligned}
\text{Präzession}&=50''{,}11\\
\text{Aberration}&=20''{,}00\\
\text{Längendifferenz Paris-Greenwich}&=9^{\mathrm m}20^{\mathrm s}{,}8.
\end{aligned}
\]

Die Werte für die Elemente (Seite 204) ergaben sich bei der Kontrollrechnung teilweise
etwas abweichend. Die Genauigkeitsgrenze, innerhalb welcher sie überhaupt zu bestimmen
sind, hängt von der Schärfe ab, mit der die Größe $x = \sin\tfrac12 g^2$ bekannt ist. In
der ersten Hypothese wurde bei der \SourceLeaf{292}Kontrollrechnung gefunden
\[
\tfrac12 g = 1^0 34' 10''{,}046 \qquad x = 0{,}000\,7501\,446
\]
und hieraus ergab sich
\[
\begin{aligned}
\text{Epoche} &= 42^0 12' 39''{,}89 \\
\text{Mittlere Bewegung} &= 823''{,}4987 \\
\text{Perihel} &= 52^0 41' 13''{,}24 \\
\varphi &= 14\ 24\ 28{,}58 \\
\text{Logarithmus der halben großen Achse} &= 0{,}422\,8958 \\
\text{Aufsteigender Knoten} &= 171^0 5' 48''{,}84 \\
\text{Neigung} &= 13^0 2' 37''{,}50.
\end{aligned}
\]

Die von \textsc{Gauss} angegebenen Werte scheinen auf der Annahme
\[
\tfrac12 g = 1^0 34' 10''{,}056 \qquad x = 0{,}000\,7501\,474
\]
zu beruhen. Die Differenzen liegen also unterhalb der zu erreichenden Genauigkeitsgrenze.
In der zweiten Hypothese wurden sie noch geringer gefunden, in der dritten ergaben sich
genau die \textsc{Gauss}schen Werte, abgesehen von den im Fehlerverzeichnis vermerkten
Stellen.

Zu Art. 156---157. Die Deklination der dritten Beobachtung der Pallas ist in
\textsc{Zach}s Monatlicher Correspondenz a.\,a.\,O. zu $-28^0 38' 8''{,}3$ angegeben. Die
Sonnenörter Seite 205, Zeile 5---7 von unten sind nicht kontrolliert. Die Nutation hat
\textsc{Gauss} vermutlich nach seinen eigenen, handschriftlich in ein Handbuch (B\,b)
eingetragenen Tafeln berechnet. Diese Vorarbeit unterscheidet sich von den in der
Monatlichen Correspondenz (s. Band VI, Seite 123) abgedruckten Tafeln nur durch andere
Werte der Konstanten. Die Nutation erwies sich als richtig
bis auf den dritten Wert ($13''{,}68$ statt $12''{,}20$). Die Aberration, mit dem Wert der
Konstanten $20''{,}00$ berechnet, erwies sich als richtig bis auf den dritten und vierten Wert ($21''{,}51$ und $1''{,}63$ statt
$23''{,}68$ und $1''{,}31$). Die Präzession hat \textsc{Gauss} mit den Werten $m = 45''{,}93$,
$n = 20''{,}02$ berechnet. Hier differiert der dritte Ort ($2''{,}55$ statt $2''{,}25$). Diese
Differenzen haben beim Abdruck nur teilweise berücksichtigt werden können, vgl. das
Fehlerverzeichnis.

Die Elemente Seite 211 fanden sich z.\,T. auch etwas abweichend, aber innerhalb der
Genauigkeitsgrenze der ganzen Rechnung, z.\,B. für das Perihel vom Knoten $323^0 44' 56''{,}64$.

Zu Art. 159. Die Beobachtungszeit des ersten Orts der Ceres ist im Briefe \textsc{Olbers}
an \textsc{Gauss} vom 8. September 1805 zu $13^{\mathrm h} 8^{\mathrm m} 47^{\mathrm s}$
angegeben. Der Reduktion sind die Längendifferenzen Bremen-Paris $= 25^{\mathrm m}
51^{\mathrm s}$, Göttingen-Paris $= 30^{\mathrm m} 21^{\mathrm s}$, Lilienthal-Paris $=
26^{\mathrm m} 14^{\mathrm s}$ zugrunde gelegt, welche in \textsc{Zach}, Tabulae motuum
Solis, Supplementum, Gotha 1804, mitgeteilt sind. Seite 213 Zeile 8---11 sind die Werte
für Schiefe der Ekliptik, Präzession und Nutation und Zeile 9 u. 8 von unten die
Sonnenörter keiner Kontrolle unterworfen worden. Seite 214 haben sich in den
Zahlenangaben für Hypothese II größere Unrichtigkeiten herausgestellt, die nur zum Teil
verbessert werden konnten. Die Werte von $\omega+\sigma$, $\log Qc\sin\omega$ und $z$
sind verbessert worden. Der darauf folgende Wert von $r'$ ergibt sich dann auch so, wie
\textsc{Gauss} ihn angibt, ist also wahrscheinlich mit dem eingesetzten richtigen Werte
von $z$ berechnet. Dagegen ist von hier ab die weitere Rechnung mit dem unrichtigen
Werte der alten Ausgabe ($z = 7^0 2' 32''{,}870$) richtig weitergeführt. Seite 215 ist die
Rechnung bis zum $\log r' = 0{,}413\,2817$ richtig. Von hier ab aber scheint
\textsc{Gauss} mit den unrichtigen Werten $\log\dfrac{n'r'}{n} = 0{,}662\,0412$ (statt
$0{,}662\,0380$) und $\log\dfrac{n'r'}{n''} = 0{,}636\,4081$ (statt $0{,}636\,4049$)
gerechnet zu haben, denen die weitere Rechnung richtig entspricht. Auch die Werte für die
wahren Anomalien, den Winkel $\varphi$ und die halbe große Achse Seite 215, Zeile 11, 10,
5, 4 v. unten ergaben sich bei der Nachrechnung etwas abweichend, sowohl mit Anwendung des
in der alten Ausgabe Art. 97 gegebenen
\SourceLeaf{293}
Wertes $x = 0{,}065\,29078$, als auch des verbesserten $x = 0{,}065\,29076$. Den von
\textsc{Gauss} angegebenen Werten würde $g = 29^0 36' 31''{,}975$ ($x = 0{,}065\,290798$)
entsprechen.

Zu Art. 171. In der ersten Beobachtung der \textls[150]{Vesta} gibt \textsc{Olbers} im Briefe an
\textsc{Gauss} vom 5. April 1807 die Rektaszension zu $183^0 52' 37''$ an. Seite 231 sind
die Sonnenörter Zeile 12---15, sowie die Ableitung der Zahlen ebenda Zeile 9---6 v. unten
aus den Beobachtungen nicht kontrolliert. Seite 235 scheinen die Werte der Elemente mit
Hilfe von $g = 23^0 34' 22''{,}34$ ($x = 0{,}041\,7240$) gerechnet zu sein. Die Zahlen Seite
235, Zeile 7 bis 2 v. unten wurden ebenfalls nicht kontrolliert.

Alle handschriftlichen Bemerkungen, welche \textsc{Gauss} in sein Handexemplar der
Theoria motus eingetragen hat, sind als Fußnoten eingefügt, und zwar Seite 20, 30, 52,
53, 56, 65, 66, 80, 117---118, 120, 144, 145, 183, 226, 242, 249, 251, 270, wobei in der
Notiz Seite 118, letzte Zeile, $\dfrac{17420}{729}$ in $\dfrac{15420}{729} =
\dfrac{5120}{243}$ und Seite 185, Zeile 7 v. unten $-3 \pm \BareRoot{(9-16\operatorname{cotang} A^2)}$ in
$-3\operatorname{cotang} A \pm \BareRoot{(9\operatorname{cotang} A^2-16)}$ verbessert worden ist. Die Entwicklung des
Ausdrucks in der Fußnote S. 270 gibt $y = 1 + \dfrac{10}{9}h + \dfrac{110}{81}hh +
\cdots$. Es ist also vermutlich (vgl. S. 118) ein Zeichenfehler begangen worden, und
sollte etwa heißen $y = \left(1+\dfrac{32}{9}h\right)^{\frac{5}{16}}$. Zu Seite 144
vgl. die Bemerkungen über \textsc{Ivory}s Methode zur Kometenbahnbestimmung im Briefe
\textsc{Gauss} an \textsc{Olbers} vom 31. Dezember 1814.

Ferner mag Folgendes bemerkt werden.

Zu Art. 32\,IV. Der Maximalfehler in der wahren Anomalie tritt für den Wert von $v$ ein,
den eine der beiden folgenden Gleichungen bestimmt.
\begin{center}
$\cos v + e\cos 2v = 0$\quad oder\quad \mbox{$\cos v = \dfrac1e\left\{
-\tfrac14 + \BareRoot{\left(\tfrac1{16}+\tfrac12 ee\right)}\right\}$}
\end{center}
Für alle auf
Seite 43 angegebenen Werte von $e$ ist $v$ ungefähr gleich~$60^0$.

Zu Art. 32\,VII. In der Hyperbel tritt der Maximalfehler für $F$ ein, wenn $\cos F +
\tfrac13\sin 2F - e\cos 2F = 0$. Die der Tabelle auf Seite 44 entsprechenden Werte von
$F$ sind

\begin{center}
\begin{tabular}{c|c}
$e$ & $F$ \\
\hline
$1{,}3$ & $14^0 46'$ \\
$1{,}2$ & $11\ 16$ \\
$1{,}1$ & $6\ 41$ \\
$1{,}05$ & $3\ 44$ \\
$1{,}01$ & $0\ 50$ \\
$1{,}001$ & $0\ 5$ \\
\end{tabular}
\end{center}

Zu Art. 36. Der Maximalfehler in der Ellipse tritt ein für $1-2ee+e\cos E = 0$, in der
Hyperbel für $ee(1+6\sin F)+e(\cos F-3\operatorname{tang} F) = 2+3\sin F$. Die zu den Seite 48
gegebenen Tabellen gehörigen Werte von $e$ sind also die folgenden

\begin{center}
\begin{minipage}[t]{.42\linewidth}\centering
Ellipse\par\smallskip
\begin{tabular}{c|c}
$E$ & $e$ \\
\hline
$10^0$ & $0{,}99494$ \\
$20$ & $0{,}98003$ \\
$30$ & $0{,}95602$ \\
$40$ & $0{,}92409$ \\
$50$ & $0{,}88583$ \\
$60$ & $0{,}83534$ \\
\end{tabular}
\end{minipage}\hfill
\begin{minipage}[t]{.42\linewidth}\centering
Hyperbel\par\smallskip
\begin{tabular}{c|c}
$F$ & $e$ \\
\hline
$10^0$ & $1{,}00508$ \\
$20$ & $1{,}02094$ \\
$30$ & $1{,}04993$ \\
$40$ & $1{,}09754$ \\
$50$ & $1{,}17671$ \\
$60$ & $1{,}32007$ \\
$70$ & $1{,}63435$ \\
\end{tabular}
\end{minipage}
\end{center}
\SourceLeaf{294}
\textsc{Gauss}’ Vorarbeiten zur Theoria motus finden sich hauptsächlich in den mit B\,b, B\,c, B\,d bezeichneten Handbüchern, die Tafeln im Handbuch B\,b S. 30, 31, 35. Zur Entstehung des Manuskriptes und für die Vorbereitungen zur Drucklegung vergleiche man den Briefwechsel mit \textsc{Olbers} (W. \textsc{Olbers}, sein Leben und seine Werke, herausgegeben von C. \textsc{Schilling}, Band II).

\begin{flushright}
\textsc{Brendel}.
\end{flushright}

\begin{center}
\rule{0.3\linewidth}{0.4pt}
\end{center}
